\documentclass[14 pt]{extarticle}

	\usepackage[french]{babel}
	\usepackage[utf8]{inputenc}  
	\usepackage[T1]{fontenc}
	\usepackage{amssymb}
	\usepackage[mathscr]{euscript}
	\usepackage{stmaryrd}
	\usepackage{amsmath}
	\usepackage{tikz}
	\usepackage[all,cmtip]{xy}
	\usepackage{amsthm}
	\usepackage{varioref}
	\usepackage{geometry}
	\geometry{a4paper}
	\usepackage{lmodern}
	\usepackage{hyperref}
	\usepackage{array}
	 \usepackage{fancyhdr}
\renewcommand{\theenumi}{\alph{enumi})}
	\pagestyle{fancy}
	\theoremstyle{plain}
	\fancyfoot[C]{} 
	\fancyhead[L]{Contrôle}
	\fancyhead[R]{19 mai 2026}\geometry{
 a4paper,
 total={170mm,257mm},
 left=20mm,
 top=20mm,
 }
	
	
	\title{Contrôle proportionnalité}
	\date{}
	\begin{document}
\emph{Toute réponse doit être soigneusement justifiée.}
 \subsection*{Exercice 1 - 4 points}
Un véhicule consomme $6$ litres d'essence en $100$ kilomètres. \begin{enumerate}
 \item Combien consomme-t-il en $350$ kilomètres ? 
 \item Quelle distance va-t-il parcourir avec $54$ litres d'essence ?
 \end{enumerate}
 \subsection*{Exercice 2 - 4 points}
 \begin{enumerate}
 \item  Calculer $30\%$ de $42$ euros. 
 \item  Calculer $15\%$ de $54$ grammes. 
 \end{enumerate}
 
 \subsection*{Exercice 3 - 4 points}
 Convertir :
 \begin{enumerate}
 \item $72$ km/h en m/s
 \item $25$ m/s en km/h.
 \end{enumerate}
 
 
 \subsection*{Exercice 4 - 4 points}
 
 Un produit qui coûtait 720 euros a augmenté de $25\%$ en cinq ans.
 \begin{enumerate}
 \item Combien coûte-t-il après cette hausse ?
 \item Le produit retrouve son prix initial l'année suivante : de combien de \% a-t-il baissé cette année-là ?
 \end{enumerate}
 \subsection*{Exercice 5 - 4 points}
Dans une classe, il y  a 12 filles et 18 garçons.
\begin{enumerate}
\item  Quel est le pourcentage de filles ?
\item On sait que cette classe contient $15$ \% des élèves du collège. Combien y a-t-il d'élèves ? 
\end{enumerate}

\subsection*{Bonus une fois fini le reste}
\begin{enumerate}
\item Un véhicule parcourt un chemin en avançant 1 h à 40 km/h, puis 1h à 80 km/h. Quelle est sa vitesse moyenne ? 
\item Un véhicule parcourt un chemin en effectuant 20 km à 40 km/h, puis 20 km à 80 km/h. Quelle est sa vitesse moyenne ? 
\end{enumerate}
  
 	\end{document}
 	\newpage
